\subsection{Event-Input Controls}

We evaluate the final reported checkpoint,
\path{outputs/night3r_mvsec_only_adapt_1k_from_best/last_event_dark3r.pth},
under real, zero, and shuffled event inputs.  All settings use the same model
weights and the same 128 fixed adjacent frame pairs.  The zero control removes
the event signal entirely.  The shuffle control preserves the event-value
distribution and event reliability statistics while breaking spatial and
temporal alignment with RGB.

\begin{table}[h]
\centering
\caption{Final Night3R checkpoint under real, zero, and shuffled event inputs.
All results use the MVSEC-only 1k final weight.  Lower AbsRel and higher
$\delta_1$ are better.}
\label{tab:event-control}
\begin{tabular}{llrrrr}
\toprule
Dataset & Event & Raw AbsRel \down & Raw $\delta_1$ \up & Median AbsRel \down & Median $\delta_1$ \up \\
\midrule
\multirow{3}{*}{DSEC} & Real & \best{0.1356} & \best{0.8373} & \best{0.0559} & \best{0.9826} \\
 & Zero & 0.3101 & 0.2689 & 0.1071 & 0.9176 \\
 & Shuffle & 0.1833 & 0.6305 & 0.0866 & 0.9188 \\
\midrule
\multirow{3}{*}{MVSEC} & Real & \best{0.2073} & \best{0.8118} & \best{0.1949} & \best{0.8016} \\
 & Zero & 0.6798 & 0.2057 & 0.4216 & 0.3514 \\
 & Shuffle & 0.3165 & 0.5754 & 0.3122 & 0.5027 \\
\bottomrule
\end{tabular}
\end{table}

Real events consistently outperform both controls.  On MVSEC, real events
reduce Raw AbsRel by 69.5\% relative to zero events and by 34.5\% relative to
shuffled events.  On DSEC, the corresponding reductions are 56.3\% and 26.0\%.
Both datasets therefore exhibit the ordering Real $<$ Shuffle $<$ Zero in Raw
AbsRel.  This shows that the final model benefits from the presence of event
measurements, while the additional degradation caused by shuffling confirms
that correct event--RGB alignment is essential.

Importantly, real and shuffled inputs have nearly identical event density,
occupancy, and mean gate values.  For MVSEC, the mean gates are 0.0203 and
0.0205 for real and shuffled events, respectively; for DSEC, both are 0.0211.
The accuracy gap therefore cannot be explained by event count, marginal event
statistics, or a larger average fusion gate.  It is caused by whether the
event residual carries correctly aligned local structure.

MVSEC Raw AbsRel (0.2073) and median-aligned AbsRel (0.1949) are close, showing
that the final checkpoint retains stable metric-scale calibration on its target
domain.  On DSEC, Raw AbsRel is 0.1356 while median-aligned AbsRel is 0.0559,
indicating accurate relative geometry but a remaining cross-domain scale
offset.  MVSEC warp PSNR changes by less than 0.1 dB across the three event
modes, so depth accuracy is a substantially more sensitive event-input metric
than this pairwise photometric diagnostic.

\paragraph{Controlled synthetic low-light depth evaluation.}
To evaluate low-light robustness independently of stored image brightness, we
darken the same 128 fixed DSEC pairs at test time.  Each sRGB image is
inverse-gamma linearized with $\gamma=2.2$, multiplied by a fixed exposure,
corrupted by signal-dependent shot noise and signal-independent read noise with
scale 0.01, clipped, and gamma encoded.  Noise is keyed by sample identifier,
so MASt3R, Dark3R, and all three Night3R controls receive pixelwise identical
darkened RGB.  The recorded event voxel is retained, zeroed, or shuffled
according to the event control.

\begin{table*}[t]
\centering
\caption{DSEC depth accuracy for the final MVSEC-only 1k checkpoint under
controlled synthetic low light.  All settings use identical darkened RGB.}
\label{tab:event-control-dsec-darken}
\begin{tabular}{clrrrr}
\toprule
Exposure & Input/model & Raw AbsRel $\downarrow$ & Raw $\delta_1$ $\uparrow$ & Median AbsRel $\downarrow$ & Median $\delta_1$ $\uparrow$ \\
\midrule
\multirow{5}{*}{0.20} & MASt3R & 0.3907 & 0.0602 & 0.1623 & 0.7494 \\
 & Dark3R & 0.5885 & 0.0002 & 0.1632 & 0.7585 \\
 & Real & 0.1611 & 0.7128 & \best{0.0710} & \best{0.9679} \\
 & Zero & \best{0.1182} & \best{0.8420} & 0.1096 & 0.8882 \\
 & Shuffle & 0.1868 & 0.6256 & 0.1410 & 0.7920 \\
\midrule
\multirow{5}{*}{0.10} & MASt3R & 0.5450 & 0.0082 & 0.2610 & 0.5263 \\
 & Dark3R & 0.6573 & 0.0000 & 0.1998 & 0.6591 \\
 & Real & 0.1727 & 0.6503 & \best{0.0796} & \best{0.9565} \\
 & Zero & \best{0.1393} & \best{0.7839} & 0.1325 & 0.8226 \\
 & Shuffle & 0.2105 & 0.5525 & 0.1684 & 0.7308 \\
\midrule
\multirow{5}{*}{0.05} & MASt3R & 0.8724 & 0.0008 & 0.4914 & 0.3486 \\
 & Dark3R & 0.8102 & 0.0000 & 0.2839 & 0.5164 \\
 & Real & \best{0.1859} & \best{0.5861} & \best{0.0934} & \best{0.9301} \\
 & Zero & 0.3479 & 0.2341 & 0.1939 & 0.7317 \\
 & Shuffle & 0.2555 & 0.4132 & 0.2028 & 0.6611 \\
\bottomrule
\end{tabular}
\end{table*}

At $e=0.05$, real events reduce Raw AbsRel by 78.7\% relative to MASt3R,
77.1\% relative to Dark3R, 46.6\% relative to zero, and 27.2\% relative to
shuffle.  They reduce Median AbsRel by 67.1\% relative to Dark3R, 51.8\%
relative to zero, and 53.9\% relative to shuffle.  At $e=0.20$
and $e=0.10$, real events remain best after median alignment, while zero events
have lower raw-scale error.  This indicates a scale-calibration offset rather
than better relative geometry for the zero control.  The monotonic MASt3R
degradation and real-over-shuffle advantage show that aligned events become
increasingly important as RGB visibility decreases.

\paragraph{DL3DV diagnostics.}
The same final checkpoint is also evaluated on 128 fixed adjacent pairs from
the full DL3DV evaluation root.  The local DL3DV format provides neither metric
depth ground truth nor usable warp supervision for this evaluator; consequently,
we report event-path diagnostics rather than reconstruction accuracy.

\begin{table}[h]
\centering
\caption{Event-path diagnostics on DL3DV.  These values verify event-control
execution and distribution preservation; they are not reconstruction accuracy
metrics.}
\label{tab:event-control-dl3dv}
\begin{tabular}{lrrr}
\toprule
Event & Gate mean & Event density & Event occupancy \\
\midrule
Real & 0.01918 & 0.5803 & 0.6115 \\
Zero & 0.00000 & 0.0000 & 0.0000 \\
Shuffle & 0.01916 & 0.5803 & 0.6115 \\
\bottomrule
\end{tabular}
\end{table}

Real and shuffled DL3DV events retain the same density and occupancy and induce
almost identical mean gates, whereas the zero input fully disables the event
path.  This equality is expected because the shuffle control uses circular
spatial and event-bin shifts: it preserves every event value and therefore
preserves all marginal count statistics exactly.  The local split has no metric
depth or usable warp-quality reference, so these diagnostics cannot by
themselves establish a depth improvement.  Its camera trajectories do, however,
permit the separate pose evaluation below.

\paragraph{Trajectory evaluation.}
We additionally recover a camera trajectory from the dense model outputs.  For
each adjacent pair, inference is run in both image orders.  The point map of the
same image expressed in camera~1 and camera~2 coordinates is registered with a
trimmed Umeyama Sim(3) fit.  Both MASt3R head directions are evaluated and the
direction with the lower normalized 3D registration residual is retained.  The
relative rotation is used directly, while relative translation is normalized
per edge because monocular pairwise reconstruction has an arbitrary scale.
Relative poses are chained for 32 frames, followed by one EVO Sim(3) alignment
per scene.  We report the mean of per-scene RMSE over six scenes.  DL3DV camera
axes are converted from OpenGL to OpenCV before evaluation.  DSEC is unavailable
because the local DSEC test export contains no camera-pose ground truth.
This table uses no test-time low-light augmentation: MVSEC is the real
\texttt{outdoor\_night2} sequence, while DL3DV records
\texttt{image\_mode=copy} and therefore uses source RGB without synthetic
darkening.

\begin{table*}[t]
\centering
\caption{Trajectory accuracy for the final MVSEC-only 1k checkpoint.  ATE and
RPE are computed over six 32-frame trajectories after one per-scene Sim(3)
alignment.  Lower is better.}
\label{tab:event-control-pose}
\begin{tabular}{llrrr}
\toprule
Dataset & Input/model & ATE (m) $\downarrow$ & RPE trans. (m) $\downarrow$ & RPE rot. ($^\circ$) $\downarrow$ \\
\midrule
\multirow{4}{*}{MVSEC} & MASt3R & \textbf{0.3705} & \textbf{0.2557} & \textbf{6.44} \\
 & Real & 0.3822 & 0.2829 & 95.77 \\
 & Zero & 0.4707 & 0.3613 & 133.19 \\
 & Shuffle & 0.4992 & 0.3086 & 111.56 \\
\midrule
\multirow{4}{*}{DL3DV} & MASt3R & \textbf{0.2570} & \textbf{0.1426} & \textbf{1.60} \\
 & Real & 0.4289 & 0.3151 & 153.71 \\
 & Zero & 0.4984 & 0.3290 & 155.16 \\
 & Shuffle & 0.3053 & 0.2516 & 139.71 \\
\midrule
DSEC & \multicolumn{4}{l}{Not available: no camera-pose ground truth in the local test export.} \\
\bottomrule
\end{tabular}
\end{table*}

On MVSEC, real events reduce ATE by 18.8\% relative to zero events and by
23.4\% relative to shuffled events.  They also reduce translational RPE by
21.7\% and 8.4\%, respectively.  Thus the pose translation metrics agree with
the depth controls on the target domain: aligned event measurements are more
useful than either removing the event signal or preserving its distribution
while breaking alignment.

This conclusion does not transfer cleanly to DL3DV.  Real events improve ATE
over zero events (0.4289 versus 0.4984), but shuffled events obtain the lower
ATE of 0.3053.  More importantly, all three Night3R inputs have very large
rotation RPE and normalized bidirectional point-map residuals above 1.0,
whereas MASt3R obtains 1.60 degrees and a residual of 0.0199 on DL3DV.  The
corresponding MVSEC residuals are 1.2295 for real events and 0.1104 for MASt3R.
Consequently, these trajectory results support a real-event advantage over the
zero/shuffle controls only on MVSEC; they do not support a claim that the final
depth-adapted checkpoint improves pose over MASt3R.  They instead reveal that
the current fine-tuning preserves useful depth ordering while degrading full
XYZ cross-view consistency and rotation recovery.

For completeness, we repeated the DL3DV trajectory protocol after deterministic
darkening at exposures 0.20/0.10/0.05.  MASt3R obtains ATE
0.2495/0.2368/0.2489 and rotation RPE 1.76/2.19/3.11 degrees.  Night3R with real
events obtains ATE 0.4502/0.4533/0.4124 and rotation RPE
156.12/156.54/156.79 degrees; zero and shuffled events remain similarly poor
and show no consistent ordering.  Explicit darkening therefore does not repair
the final checkpoint's cross-view XYZ inconsistency.  This negative result is
retained as a diagnostic and is not used to claim low-light pose improvement.
